% Elliot Alladaye — CV (English)
% Self-contained: compiles with pdflatex using only base TeXLive packages.
\documentclass[10pt,a4paper]{article}

\usepackage[T1]{fontenc}
\usepackage[utf8]{inputenc}
\usepackage[provide=*,french]{babel}
\usepackage[top=1.1cm,bottom=0.9cm,left=1.4cm,right=1.4cm]{geometry}
\usepackage{titlesec}
\usepackage{enumitem}
\usepackage{xcolor}
\usepackage[hidelinks]{hyperref}
\usepackage{array}

\definecolor{accent}{HTML}{1F4E79}

\setlength{\parindent}{0pt}
\linespread{0.95}
\pagestyle{empty}

\titleformat{\section}
  {\Large\bfseries\color{accent}}{}{0pt}{}[\vspace{1pt}\hrule height 0.8pt\vspace{4pt}]
\titlespacing*{\section}{0pt}{3pt}{2pt}

\setlist[itemize]{leftmargin=1.4em,itemsep=0.5pt,topsep=1pt,parsep=0pt}

\newcommand{\entry}[2]{%
  \vspace{1pt}\noindent\begin{tabular*}{\textwidth}{@{}l@{\extracolsep{\fill}}r@{}}
  \textbf{#1} & #2 \\
  \end{tabular*}\par}

\newcommand{\subentry}[1]{\noindent\textit{#1}\par\vspace{1pt}}

\newcommand{\skillrow}[2]{\noindent\textbf{#1:}~#2\par}

\begin{document}

% ===================== HEADER =====================
\begin{center}
  {\Huge\bfseries Elliot Alladaye}\\[3pt]
  {\large\color{accent} Ingénieur Backend \& IA}\\[5pt]
  \small
  \href{mailto:alladayeelliot@gmail.com}{alladayeelliot@gmail.com} ~$\bullet$~
  +33\,6\,XX\,XX\,XX\,XX ~$\bullet$~
  \href{https://www.linkedin.com/in/elliot-alladaye}{linkedin.com/in/elliot-alladaye}\\[2pt]
  Nord de la France / Belgique ~$\bullet$~ \textbf{Disponible en CDI/CDD à partir de septembre 2026}
\end{center}
\vspace{2pt}

% ===================== PROFILE =====================
\section{Profil}
Ingénieur Backend \& IA avec 2 ans d'expérience dans la mise en production d'outils d'IA pour les jumeaux numériques, sur tout le cycle d'intégration de l'IA : orchestration de LLM, embeddings, RAG et détection d'anomalies. Je conçois des systèmes fiables et extensibles, avec une forte exigence de confidentialité et de souveraineté des données. \textbf{MSc en Intelligence Artificielle} (University of Kent, obtenu en sept. 2026) ; diplôme d'ingénieur (Epitech, 5 ans) obtenu en 2027.

% ===================== EXPERIENCE =====================
\section{Expérience}

\entry{Ingénieur Backend \& IA --- Stereograph}{Juil. 2024 -- Août 2025}
\subentry{Temps partiel puis temps plein}
\begin{itemize}
  \item Industrialisation d'un proof-of-concept d'IA en produit prêt pour la production, puis maintenance de ce produit et de plusieurs services Python.
  \item Conception d'un \textbf{système RAG local-first} sur les documents clients, gardant les données sensibles hors des API tierces pour garantir confidentialité et souveraineté des données.
  \item Mise en production d'une détection d'anomalies sur objets 3D signalant automatiquement les valeurs aberrantes, réduisant le contrôle qualité manuel de \textasciitilde X\%.
  \item Migration de la journalisation vers \textbf{Datadog sur 4 services Python}, réduisant le temps de diagnostic des incidents.
\end{itemize}

\entry{Stagiaire Full-stack \& IA --- Stereograph}{Juil. -- Déc. 2023}
\subentry{Stage}
\begin{itemize}
  \item Développement d'un chatbot (Python + ReactJS) intégré au logiciel de jumeau numérique BIM de Stereograph, permettant l'interrogation des données de bâtiment en langage naturel.
  \item Création d'une interface textuelle pour modifier la vue 3D d'un jumeau numérique.
  \item Validation de la faisabilité technique de l'intégration de l'IA, débouchant sur une proposition d'embauche pour poursuivre le produit.
\end{itemize}

% ===================== PROJECTS =====================
\section{Projets}

\entry{Saver --- Tech Lead \& seul développeur Backend}{2023 -- présent}
\subentry{Plateforme de favoris augmentée par l'IA, de niveau production (Python 3.13, FastAPI)}
\begin{itemize}
  \item Responsable de \textbf{l'ensemble du backend} (architecture, API, base de données, pipeline IA, sécurité, observabilité) d'un service qui ingère des URL, en extrait le contenu et le classe automatiquement.
  \item Conception d'un backend asynchrone FastAPI + PostgreSQL/pgvector avec \textbf{recherche hybride} (cosinus sémantique pondéré + classement plein-texte) sur un pipeline Mistral AI.
  \item Couche d'authentification durcie : hachage Argon2, jetons JWT et jetons de rafraîchissement rotatifs avec révocation automatique de la famille de jetons en cas de vol détecté.
  \item Récupération d'URL \textbf{protégée contre les SSRF} (épinglage DNS, validation par redirection) sur 15 extracteurs spécifiques, avec repli par impersonation de navigateur.
  \item Observabilité de bout en bout : logs structurés via Fluentd $\rightarrow$ OpenTelemetry $\rightarrow$ Loki $\rightarrow$ Grafana, 5 dashboards en code ; \texttt{mypy} strict et tests asynchrones complets.
\end{itemize}

\entry{AREA --- Moteur d'automatisation (type Zapier)}{}
\begin{itemize}
  \item Développement d'un moteur d'automatisation en microservices (Node.js + TypeScript), séparant l'API des workers d'arrière-plan pour une meilleure modularité.
  \item Implémentation d'une file de tâches fiable (RabbitMQ) avec patterns de retry, garantissant l'exécution des workflows malgré les défaillances transitoires.
\end{itemize}

% ===================== EDUCATION =====================
\section{Formation}

\entry{MSc en Intelligence Artificielle --- University of Kent}{2025 -- 2026}
\subentry{Diplôme obtenu en septembre 2026 --- année de double diplôme à l'international du cursus Epitech.}

\entry{Diplôme d'ingénieur (niveau Master), Génie logiciel --- Epitech}{2022 -- 2027}
\subentry{Grande École française, cursus de 5 ans ; année à l'étranger à l'University of Kent. Diplôme obtenu à l'été 2027.}

% ===================== SKILLS =====================
\section{Compétences}
\skillrow{Langages}{Python, TypeScript, C/C++}
\skillrow{IA / ML}{Orchestration de LLM, embeddings \& RAG, détection d'anomalies ; NumPy, Pandas, scikit-learn}
\skillrow{Backend \& Infra}{FastAPI, PostgreSQL/pgvector, SQLAlchemy, API REST, ReactJS, Docker, Datadog, Grafana, Loki}
\skillrow{Langues}{Français (langue maternelle), Anglais (bilingue)}

\end{document}
